\documentclass[10pt]{article}
\usepackage[margin=1in]{geometry}
\usepackage{booktabs}
\usepackage{graphicx}
\usepackage{amsmath,amssymb}
\usepackage{lmodern}
\usepackage[T1]{fontenc}
\usepackage{microtype}
\usepackage{longtable}
\usepackage[hidelinks]{hyperref}

% Shared macros. Names and numbers that appear more than once live here so that a
% value cannot drift between two mentions of it.

\newcommand{\pns}{prior-normalised skill}
\newcommand{\pss}{patient-specific skill}
\newcommand{\Pns}{Prior-normalised skill}
\newcommand{\Pss}{Patient-specific skill}

% The two objects previously both called "fitted template". Never write either
% phrase inline; renaming one edits these four lines and nothing else.
\newcommand{\maskref}{mask-fitted reference}        % fit to validation lesion masks; the content-free ceiling
\newcommand{\Maskref}{Mask-fitted reference}
\newcommand{\attden}{attention-fitted denominator}  % reference_baselines.fitted_template; fit to the arm's own attention
\newcommand{\Attden}{Attention-fitted denominator}

% Numbers quoted more than once, from the confirmatory artefacts.
\newcommand{\maskrefInplane}{0.8302}   % masks:inplane flat AUC
\newcommand{\maskrefSep}{0.8543}       % masks:separable flat AUC
\newcommand{\attdenMedian}{0.4208}     % median attention:inplane over the 40 H5 dumps
\newcommand{\probeDen}{0.8358}         % the probe's attention-fitted denominator
\newcommand{\centrePriorSlice}{0.6427} % model-free centre prior, slice AUC, \condA{}
\newcommand{\armSlot}{slot arm}

\newcommand{\preghash}{\texttt{4fd6e801157ecef5}}
\newcommand{\repourl}{https://github.com/BhaveshThapar/SlotMIL}

\newcommand{\armNG}{Normal Guidance}
\newcommand{\armGA}{gated ABMIL}
\newcommand{\condA}{\texttt{nodule\_present}}
\newcommand{\condB}{\texttt{malignancy}}
\newcommand{\condC}{\texttt{balanced\_presence}}
\newcommand{\condD}{\texttt{mosmed\_severity}}

% A verdict, typeset so a FAIL is never quieter than a PASS.
\newcommand{\PASS}{\textsc{pass}}
\newcommand{\FAIL}{\textbf{\textsc{fail}}}
\newcommand{\RPT}{\textsc{reported}}


\title{Companion to ``Position priors dominate instance-level localisation
metrics for weakly-supervised multiple-instance learning in volumetric CT''}
\author{
  Bhavesh Thapar\\
  \texttt{bthapar@terpmail.umd.edu}
  \and
  Mohammad Nayeem Teli\\
  \texttt{nayeem@umd.edu}
}
\date{}

\begin{document}
\maketitle

This companion carries the apparatus the main paper summarises in one sentence:
the full falsifier table, the diff between what was declared and what was
reported, and the per-condition detail. Everything here is regenerable from the
repository at configuration hash \preghash{}.

% ==========================================================================
\section{The full falsifier table}

Statements and thresholds are quoted from the frozen configuration; observed
values are quoted from the verdict artefact for \condA{}, the condition carrying
the hypothesis family.

\begin{center}\small
\begin{longtable}{@{}p{0.5cm}p{5.4cm}p{1.9cm}p{5.4cm}p{1.5cm}@{}}
\toprule
& Pre-registered statement & Bar & Observed & Verdict \\
\midrule
\endhead
H1 & For every arm, |flat instance AUC - within-slice instance AUC| \textless{} 0.02: the reported 3D metric is an in-plane metric wearing 3D clothing. & 0.02 & 1/8 arms have |flat - within\_slice| \textgreater{}= 0.02; a majority is 5 & PASS \\
\addlinespace[2pt]
H2 & The model-free centre prior's slice AUC exceeds every trained arm's slice AUC -- attention MIL is worse than geometry on the axial axis. & -- & 0/8 arms match or beat the centre prior's slice AUC of 0.6427 & PASS \\
\addlinespace[2pt]
H3 & Over \textgreater{}=30 untrained inits the 95th-percentile instance AUC exceeds 0.70, i.e. 0.5 is not the floor. & 0.7 & 95th-percentile untrained instance AUC 0.7666 vs 0.7 & PASS \\
\addlinespace[2pt]
H4 & Every arm's patient-specific skill is \textless{} 0.15, and the median across arms is \textless{} 0.10. & per\_arm 0.15, median 0.1 & every arm below 0.15 and median -0.0466 below 0.1 & PASS \\
\addlinespace[2pt]
H5 & No arm's prior-normalised skill, taken as the mean over that arm's float32 seeds, exceeds 0.30. & 0.3 & 0/8 arms exceed 0.3 on the mean over seeds & PASS \\
\addlinespace[2pt]
H6 & Normal Guidance raises slice-level AUC over its base arm but raises patient-specific skill by less than 0.02. & 0.02 & patient-specific skill moved +0.1509 against a falsifier at \textgreater{}= 0.02 & \textbf{FAIL} \\
\addlinespace[2pt]
H7 & The supervised patch probe's prior-normalised skill exceeds 0.50 while every content-free baseline's is below 0.05. & probe 0.5, content\_free 0.05 & probe 0.7305 \textgreater{} 0.5 and every computed content-free member \textless{} 0.05 & PASS \\
\addlinespace[2pt]
H8 & In-lung fitted-template AUC exceeds 0.65. & 0.65 & in\_lung\_auc 0.4884 is below 0.65 (two-sided, no falsifier) & REPORTED \\
\addlinespace[2pt]
H9 & MosMed's fitted-template AUC is lower than LIDC's, with correspondingly higher prior-normalised skill -- prior dominance tracks the positional stereo… & -- & 1/50 matched tags order MosMed's template AUC below LIDC's with higher skill & \textbf{FAIL} \\
\addlinespace[2pt]
H10 & Per seed, the paired per-bag difference between the mask-fitted separable template and the trained arm on the flat axis does not favour the trained a… & -- & trained\_wins on 0/5 seeds; a majority is 3 & PASS \\
\addlinespace[2pt]
\bottomrule
\end{longtable}

\end{center}

\noindent Where the frozen prose says ``fitted template'' it names the
\attden{} --- \texttt{reference\_baselines.fitted\_template}, fit to the arm's
own attention (H5, H8, H9) --- never the \maskref{} that sets the content-free
ceiling; H10 names its mask-fitted separable member explicitly. The main
paper's terminology paragraph (\S3) separates the two.

\section{Verdicts across the three LIDC conditions}

\begin{center}
\begin{tabular}{lcccccccccc}
\toprule
Condition & H1 & H2 & H3 & H4 & H5 & H6 & H7 & H8 & H9 & H10 \\
\midrule
nodule\_present$^\dagger$ & PASS & PASS & PASS & PASS & PASS & \textbf{FAIL} & PASS & RPT & \textbf{FAIL} & PASS \\
malignancy & PASS & PASS & PASS & PASS & PASS & \textbf{FAIL} & PASS & RPT & \textbf{FAIL} & PASS \\
balanced\_presence & PASS & PASS & PASS & PASS & PASS & \textbf{FAIL} & \textbf{FAIL} & RPT & \textbf{FAIL} & \textbf{FAIL} \\
\bottomrule
\end{tabular}

\end{center}

\noindent Two entries differ from the family-carrying condition, both on
\condC{} and both attributable to its $38$ test bags: H7 fails on the probe half
($0.4772$ against $0.50$, \S\ref{sec:h7}) and H10 returns
\emph{indistinguishable} more often than the interval test can resolve.
\condD{} does not appear because it runs a strict subset of the analysis --- the
content-free template family only --- and supports H9 alone.

% ==========================================================================
\section{Declared versus reported: the diff}

\subsection{Four conditions are declared; three appear in the verdict table}

The frozen configuration declares four conditions: \condA{}, \condB{}, \condC{}
and \condD{}. Three are LIDC and receive a full verdict table. \condD{} runs the
content-free template family only and supports H9, because $50$ of $1110$ MosMed
studies carry masks and all are at severity CT-1, so re-splitting would leave too
few annotated scans in test to estimate anything. Where the main paper says
``three conditions'' it means the three LIDC conditions.

\textbf{No declared condition was dropped.} \condB{} was declared at the freeze,
received a seed-2027 patient-grouped confirmatory split on 2026-08-15, ran all
ten arms at five seeds, and reached a verdict on every hypothesis. Its stated
purpose was recorded at declaration: it has a supervised patch-probe ceiling of
$0.6516$ against \condA{}'s $0.9102$, so it tests the label rather than the
method, and it was declared a consistency check for that reason before any
result existed.

\subsection{Units taken but not declared}

Three hypotheses state a bound without a unit of aggregation. The choices, and
when each was made:

\begin{itemize}\itemsep2pt
\item \textbf{H1, mean over seeds.} A per-arm bound on
$|{\text{flat}}-{\text{within-slice}}|$ with no unit over the five seeds. Taken
as the arithmetic mean over seeds. H5's own declared rationale makes exactly this
argument for the identical hole: a per-seed maximum takes five draws at the
threshold instead of one.
\item \textbf{H4, mean over seeds.} Same shape, same reasoning.
\item \textbf{H9, every matched tag.} An ordering with no aggregation over arms
and seeds. Taken as the strictest available reading --- one matched pair failing
to order fails the hypothesis --- which is the reading least favourable to the
paper's own prediction.
\end{itemize}

All three are recorded in the verdict artefact under
\texttt{undeclared\_units\_taken} rather than left implicit, and each was chosen
before the numbers it aggregates were read. H7's aggregation unit was the fourth
such hole; it was promoted into the configuration itself on 2026-08-17 as
\texttt{content\_free\_unit} rather than recorded as taken.

\subsection{H5's denominator, and the recomputation}
\label{sec:h5}

The pre-registered denominator for \pns{} is the \attden{}: an in-plane
template fit on validation to \emph{the arm's own attention}. Fit to the
supervised probe's output that identical rule scores $\probeDen$; fit to a MIL
arm's attention its median over the $40$ scored dumps on \condA{} is
$\attdenMedian$, and it sits below $0.5$ for $24$ of them. Since $s = (A - A_t)/(1 - A_t)$, a sub-chance $A_t$
inflates the denominator above $1$ and shrinks $s$ --- and H5 passes by staying
\emph{below} a bar, so the arithmetic makes it easier to pass.

We therefore recompute the identical statistic with $A_t$ floored at $0.5$:
floor per tag, mean over the arm's seeds, maximum over the scored arm set, bar
read from the frozen configuration. This is \textbf{not} a pre-registered
analysis, it does not change H5's verdict, and it did not require amending the
configuration.

\begin{center}
\begin{tabular}{lccc}
\toprule
Condition & max skill, as published & denominator floored at 0.5 & bar \\
\midrule
nodule\_present & +0.1990 & +0.1990 & 0.3 \\
malignancy & +0.2812 & +0.2304 & 0.3 \\
balanced\_presence & +0.0885 & +0.0220 & 0.3 \\
\bottomrule
\end{tabular}

\end{center}

\noindent The verdict does not move in any condition, and flooring lowers the
maximum wherever it binds. Regenerate with
\texttt{scripts/h5\_floored\_denominator.py}.

\subsection{H7's replication was added after its single-draw values were seen}
\label{sec:h7}

Stated plainly because it matters for how the gate is read. H7's two stochastic
content-free members were originally computed once per tag: the random generator
was rebuilt per tag with a fixed seed, so every tag in a condition shared one
realisation --- the roll offsets were literally identical across all tags --- and
the spread across tags was arm-to-arm variation with the draw held constant. No
draw-to-draw variance was estimable anywhere.

Under that construction, and under the pre-registered maximum-over-tags
aggregation, the gate returned $0.0784$ and \textbf{H7 failed}. Draw $0$ of the
replication, seeded identically, gives $0.0829$ on the current arm set. The pass
does not hold on the single-draw form and we do not claim it does.

What the amendment of 2026-08-17 changed, and what it deliberately did not:

\begin{itemize}\itemsep2pt
\item Each member is now drawn $30$ times per tag and each tag's reported skill
is the mean over its draws. Thirty because the protocol already committed to
``$\geq 30$'' for the one other place the document makes a distributional claim;
reusing a number the document already stood behind is a better justification than
a fresh one chosen while looking at a result.
\item \textbf{The aggregation rule is unchanged.} The gate still takes the
maximum over tags --- the stricter choice --- because ``every content-free
baseline is below $0.05$'' is a statement about the set, and an average over tags
would let a well-behaved arm hide a badly-behaved one.
\item Draw $0$ keeps the original seed, so the published single realisation stays
locatable inside its own distribution and is emitted per tag.
\item \textbf{Results already seen: yes.} The single-draw numbers had been seen,
including the maximum that put H7 at \FAIL{}, and the full per-member table under
mean, median and max was computed for all three conditions before the amendment
was drafted --- it showed the gate clearing under mean or median in every
condition and failing under max in every condition. The amendment was nonetheless
written to leave the maximum in force, and H7's outcome after replication was not
predicted in it.
\end{itemize}

The decisive quantity is the replicated spread. Per-tag draw-to-draw standard
deviations reach $0.0673$ (entropy-matched random) and $0.0483$ (roll
permutation) on \condA{}, against the $0.05$ threshold a single draw was being
compared to. A maximum over one draw from a distribution whose spread exceeds the
threshold is not an unlucky measurement; it is not a measurement. The verdict
artefact carries \texttt{aggregation\_sensitivity} under all three aggregations
so the choice is checkable rather than assertable.

\subsection{Two H9 counts that measure different things}

The verdict artefact's \texttt{reason} reports $1$ of $50$ matched tags and its
\texttt{power.n\_ordering} reports $2$. Both are correct and they are different
quantities: \texttt{reason} counts tags clearing the \emph{full} falsifier
(ordering \emph{and} non-overlapping intervals), \texttt{n\_ordering} counts tags
ordering on the point estimate alone regardless of intervals.

\subsection{TransMIL was promoted after measuring}

TransMIL carried a drop rule --- if it was not training by 2026-09-22 it was cut
and the paper reported the arm count achieved. The promotion was made
\emph{after} measuring rather than before, because promoting first would have
re-hashed the configuration and superseded every stamp on disk to answer one
question: does an architecture designed for whole-slide bags train on LIDC's
${\sim}43.8$k-instance bags? It does, at $16.22$ of $23.55$ GiB peak. The rule did
not fire and the paper reports ten arms rather than nine. The smoke job that
decided it ran on the discovery split under an exploratory role; no hypothesis
reads its number.

\section{Reproducing}

Every analysis driver is a standalone CLI reading cached attention dumps, needs
no GPU, and writes a stamped artefact. The scorer recomputes nothing --- it reads
only stamped JSON --- so all ten verdicts re-derive in seconds:

\begin{quote}\small\ttfamily
python scripts/prereg\_verdict.py --dir runs/nulls\_nodule\_present\_confirmatory \ldots
\end{quote}

\noindent Two sensitivity analyses added for the paper's H2 interval and
Table~3 follow the same contract:
\texttt{scripts/h2\_paired\_slice.py} (the paired patient-clustered interval
on the centre-prior margin) and \texttt{scripts/arm\_in\_lung.py} (the arms'
own raw-versus-in-lung AUC), both stamped
\texttt{sensitivity\_outside\_confirmatory\_family} and both refusing rather
than proceeding when their recomputed values disagree with the stored
confirmatory artefacts. \texttt{scripts/harvey\_bridge.py} and
\texttt{scripts/confound\_atlas.py} carry the exploratory companion analyses
and are stamped accordingly.

\noindent Figures and tables in both documents regenerate from JSON alone with
\texttt{scripts/make\_paper\_figures.py --out paper/figures --tables
paper/tables}, which also writes \texttt{paper/figures/figure\_data.json}: every
plotted number, keyed by figure and series, with its source artefact and that
artefact's configuration hash and commit. A figure that disagrees with that file
is a bug in the figure.

The chain itself verifies with \texttt{python scripts/prereg\_freeze.py --check},
which confirms the document's hash matches the committed configuration, that
every declared split still hashes as recorded, and that no artefact on disk
carries a configuration hash absent from the amendment chain.

\end{document}
